\bookchapter{第二十三章附：GitHub入門指南——從程式碼託管到協作平臺}

\chapterintro{上一章我們講了“程式設計思維”，而真正要把程式設計思維落地、和他人一起做專案，就繞不開一個名字——GitHub。無論你是計算機系學生要參加比賽、想儲存自己的程式碼作品，還是單純想看看全球開發者在折騰什麼，GitHub都是你必須認識的平臺。這一節我們從最基礎的概念講起，幫你建立一個清晰的認知。}

\booksection{一、GitHub到底是什麼？}

簡單來說，GitHub是一個基於Git的\textbf{程式碼託管與協作平臺}。它於2008年上線，2018年被微軟收購，至今仍是全球最大的開原始碼託管平臺，託管著數億個倉庫，覆蓋幾乎所有主流程式語言。

如果用一個類比來理解：\textbf{Git是你電腦上的日記本，GitHub就是一個公開的部落格平臺}——你可以釋出自己的作品，關注別人的專案，參與討論和協作。它不僅是“程式設計師的後臺倉庫”，更像是全球開發者共建的一座“技術城市”——專案在這裡誕生、迭代、被討論、被複用。

對普通人來說，GitHub也並非遙不可及。很多非程式設計師會用它來託管筆記、電子書、個人部落格，甚至用來備份配置檔案。只要你有一點點“想把自己的東西整理好並分享出去”的需求，GitHub都值得你瞭解。

\booksection{二、Git和GitHub的區別（關鍵概念！）}

很多人第一次接觸時會混淆這兩個詞，這裡幫你梳理清楚：

\begin{center}
\begin{tabularx}{\linewidth}{|X|X|X|}
\hline
\textbf{對比項} & \textbf{Git} & \textbf{GitHub} \tabularnewline
\hline
本質 & 版本控制工具（軟體） & 程式碼託管平臺（網站） \tabularnewline
執行位置 & 本地電腦上 & 雲端（瀏覽器訪問） \tabularnewline
核心職責 & 記錄檔案的每次變更（怎麼改的、誰改的、何時改的） & 把Git倉庫放到雲端，方便儲存、展示、協作 \tabularnewline
是否需要網路 & 不需要，離線可用 & 需要聯網訪問 \tabularnewline
歸屬 & 開源軟體，由Linus Torvalds（Linux之父）建立 & 微軟旗下產品 \tabularnewline
同類競品 & SVN、Mercurial & GitLab、Gitee（碼雲）、Bitbucket \tabularnewline
\hline
\end{tabularx}
\end{center}

簡單說：\textbf{Git管“怎麼記錄變更”，GitHub管“變更放在哪、怎麼協作”}。Git是工具，在本地執行；GitHub是平臺，把Git倉庫放到雲端，方便協作和展示。你可以只用Git不用GitHub（比如自己寫程式碼做版本管理），但用GitHub一定要用Git（因為GitHub就是基於Git的）。

\booksection{三、GitHub的核心概念與功能}

\booksubsection{基礎概念}

在動手之前，先認識幾個高頻詞彙，否則看GitHub的介面會一頭霧水：

\begin{itemize}
  \item \textbf{倉庫（Repository，簡稱Repo）}：一個專案對應一個倉庫，裡面存放著這個專案的所有程式碼檔案、歷史版本和檔案。可以理解為一個"專案資料夾"。
  \item \textbf{提交（Commit）}：對程式碼的一次修改記錄。每次Commit都附帶說明資訊，記錄"這次改了什麼、為什麼改"。這是Git/GitHub能"時光倒流"的基礎。
  \item \textbf{分支（Branch）}：從主線上分出來的獨立開發線。比如你想加個新功能，可以開一個分支單獨開發，開發完再合併回主線，不影響其他人。
  \item \textbf{克隆（Clone）}：把雲端倉庫完整複製一份到本地電腦，之後就可以在本地修改程式碼。
  \item \textbf{推送（Push）/拉取（Pull）}：Push是把原生程式碼上傳到GitHub，Pull是從GitHub下載最新程式碼到本地。這是團隊協作時同步程式碼的兩個核心動作。
  \item \textbf{合併請求（Pull Request，簡稱PR）}：當你在一個分支上改完程式碼，想合併到主分支時，發起一個PR，讓團隊成員審查程式碼後再合併。這是GitHub協作的核心機制。
  \item \textbf{議題（Issue）}：用來記錄bug、功能需求、討論話題。每個Issue像一個討論帖，可以指派負責人、打標籤。
  \item \textbf{Fork}：把別人的倉庫複製一份到自己帳號下，你可以在自己的副本上隨意修改，不影響原倉庫。改好後可以給原作者發PR，請求合併你的改進。
\end{itemize}

\booksubsection{核心功能}

GitHub提供的不只是程式碼儲存，還包括：

\begin{itemize}
  \item \textbf{版本控制}：記錄每次程式碼更改的歷史，改壞程式碼隨時回到上一個正常版本。
  \item \textbf{團隊協作}：支援多使用者對同一專案的並行開發，多人同時開發不會覆蓋、不會衝突。
  \item \textbf{問題跟蹤}：用Issue記錄bug和功能請求，方便團隊跟進。
  \item \textbf{專案管理}：整合了待辦事項、專案看板（Projects）等管理工具，相當於內建了一個輕量級Trello。
  \item \textbf{自動化}：透過\textbf{GitHub Actions}實現程式碼提交後自動編譯、測試、打包、部署，是現代"CI/CD（持續整合/持續部署）"的核心工具。
  \item \textbf{檔案管理}：README介紹專案，Wiki存放詳細架構和開發規範，Pages甚至可以直接託管靜態網站。
  \item \textbf{程式碼託管與展示}：每個倉庫都有漂亮的網頁展示，支援語法高亮，是你技術作品的"門面"。
  \item \textbf{開源社群}：你可以Star（收藏）喜歡的專案、Watch（關注）專案動態、Follow（關注）開發者，參與全球開源協作。
\end{itemize}

\booksection{四、註冊與上手：五分鐘建立你的第一個倉庫}

\booksubsection{註冊帳號}

\begin{enumerate}
  \item 訪問\href{https://github.com}{github.com}，點選"Sign up"。
  \item 填寫郵箱、設定密碼和使用者名稱（使用者名稱會出現在你的個人主頁地址裡，建議用英文，比如\texttt{github.\allowbreak com/\allowbreak yourname}）。
  \item 完成郵箱驗證和人機驗證。
  \item 選擇免費方案（Free）即可，對個人使用者完全夠用。
\end{enumerate}

\booksubsection{建立第一個倉庫}

\begin{enumerate}
  \item 登入後點選右上角"+"號，選"New repository"。
  \item 填寫倉庫名（Repository name），比如\texttt{my-\allowbreak first-\allowbreak project}。
  \item 寫一段簡短描述（Description），讓別人知道這個倉庫是幹嘛的。
  \item 選擇公開（Public）或私有（Private）。公開倉庫所有人可見，私有隻有自己和協作者可見。免費帳號也支援無限私有倉庫。
  \item 勾選"Add a README file"——README是專案的"說明書"，每個倉庫都建議有。
  \item 點選"Create repository"，你的第一個GitHub倉庫就誕生了！
\end{enumerate}

\begin{tcbinfo}{小知識：README是什麼}
README.md是一個Markdown格式的檔案，是別人開啟你倉庫時第一眼看到的內容。一個好的README通常包含：專案簡介、安裝方法、使用示例、許可證、聯絡方式。寫好README是程式設計師的基本素養，也是專案能否被他人使用的第一道門檻。

\end{tcbinfo}

\booksection{五、基本工作流：從本地到雲端}

理解GitHub的工作流，核心是記住這個迴圈：\textbf{本地修改 暫存 提交 推送 雲端更新}。下面用大白話解釋每一步：

\begin{enumerate}
  \item \textbf{修改檔案}：在本地電腦上用編輯器寫程式碼或改檔案。
  \item \textbf{git add}：把修改的檔案加入"暫存區"，相當於告訴Git"這些改動我準備記錄下來"。
  \item \textbf{git commit}：把暫存區的改動正式記錄成一個"提交"，附上一句說明，比如"修復登入bug"。這一步只在本地完成。
  \item \textbf{git push}：把本地的提交上傳到GitHub雲端，這樣別人就能看到了。
  \item \textbf{git pull}：在開始工作前，先從雲端拉取最新程式碼，確保本地和團隊其他人保持同步。
\end{enumerate}

對新手來說，記住這個口訣：\textbf{“先pull後改，改完add，commit帶說明，最後push上去”}。養成這個習慣，團隊協作就不會亂。

\booksection{六、團隊協作流程：Pull Request是怎麼工作的}

GitHub最強大的地方在於它的協作機制。一個典型的團隊協作流程是這樣的：

\begin{enumerate}
  \item \textbf{開分支}：從主分支（通常叫main或master）開一個新分支，比如\texttt{feature/\allowbreak login}，在這個分支上開發新功能。
  \item \textbf{開發並提交}：在新分支上寫程式碼，多次commit記錄進度。
  \item \textbf{推送到GitHub}：把分支push到雲端倉庫。
  \item \textbf{發起Pull Request}：在GitHub網頁上發起PR，請求把你的分支合併到主分支。寫清楚這個PR做了什麼改動、為什麼改。
  \item \textbf{程式碼審查（Code Review）}：團隊成員檢視你的程式碼，提出意見或建議修改。你可以根據反饋繼續commit，PR會自動更新。
  \item \textbf{合併}：審查透過後，點選"Merge"把程式碼合併到主分支。你的貢獻就正式進入專案了！
\end{enumerate}

這套流程保證了\textbf{任何程式碼進入主分支前都經過審查}，大大降低了出錯機率。開源專案也是用這套機制接受全世界陌生人的貢獻——你完全可以在某個知名開源專案上提交PR，被採納後你就是這個專案的貢獻者了。

\booksection{七、GitHub的典型應用場景}

\booksubsection{學生參加比賽與課程專案}

計算機系學生參加ACM、駭客松、畢業設計，幾乎都離不開GitHub。它既能作為團隊協作平臺，又能作為作品展示視窗——比賽評委可以直接看你的倉庫，瞭解程式碼質量和協作過程。很多公司招聘時也會要求附上GitHub主頁，作為技術能力的證明。

\booksubsection{個人知識庫與筆記託管}

非程式設計師也能用GitHub。很多人用它託管學習筆記、技術部落格（配合GitHub Pages可以免費建站）、電子書。Markdown格式寫作 + Git版本管理 + 雲端備份，三合一。

\booksubsection{參與開源專案}

GitHub是全球開源運動的核心陣地。你可以Fork喜歡的專案，修改後發PR回饋社群；也可以單純Star收藏好專案，跟蹤更新。參與開源是提升技術、積累聲譽的絕佳途徑。

\booksubsection{自動化部署與CI/CD}

藉助GitHub Actions，可以實現“程式碼一push，自動測試+自動部署到伺服器”的全自動流程。很多個人部落格、檔案站點都是這麼搭建的，寫完文章push一下，網站就自動更新了。

\booksubsection{程式碼備份與多裝置同步}

把程式碼倉庫當作雲盤用，換電腦、多裝置辦公時，git pull一下就能拿到最新程式碼，再也不用隨身碟拷來拷去。

\booksection{八、給新手的幾點建議}

\begin{itemize}
  \item \textbf{提交說明要寫清楚}：別寫"update"這種廢話說明。好的commit message寫"修復使用者登入後跳轉錯誤頁面的問題"，半年後回看也知道當時幹了啥。
  \item \textbf{勤提交，小步快跑}：每完成一個小功能就commit一次，而不是寫了一整天才提交一次。這樣出問題更容易定位。
  \item \textbf{用分支開發，別直接改主分支}：哪怕是個人專案，養成開分支的習慣，未來團隊協作時就不會手忙腳亂。
  \item \textbf{善用.gitignore}：這個檔案用來告訴Git哪些檔案不要納入版本管理（比如密碼配置、編譯產物、依賴資料夾），避免把垃圾檔案傳上去。
  \item \textbf{多看優秀開源專案}：逛逛GitHub Trending，看看高星專案的程式碼結構、README寫法、協作流程，是提升最快的免費方式。
  \item \textbf{申請GitHub學生包}：學生認證後可免費獲得大量付費工具和服務的額度（Copilot、域名、雲服務等），價值上千美元。
\end{itemize}

\begin{tcbwarn}{安全提醒}
\begin{itemize}
  \item 千萬不要把密碼、API金鑰、資料庫連線串等敏感資訊提交到GitHub，尤其是公開倉庫——一旦洩露，會被爬蟲秒抓，後果嚴重。
  \item 不小心提交了敏感資訊，光刪除檔案不夠（歷史記錄裡還有），要用\texttt{git filter-\allowbreak branch}或專業工具清理歷史，並立即更換洩露的金鑰。
  \item 私有倉庫也不是絕對安全，敏感資訊最好用環境變數管理，不程序式碼倉庫。
\end{itemize}

\end{tcbwarn}

\begin{tcbinfo}{國內訪問GitHub的注意事項}
GitHub在國內訪問有時不穩定，可能出現打不開、克隆慢的情況。可以嘗試：(1) 使用國內映象站（如GitClone、ghproxy）加速克隆；(2) 配置Hosts或使用代理；(3) 如果主要面向國內協作，也可以考慮國產替代品\textbf{Gitee（碼雲）}，功能類似且訪問速度快。很多團隊會同時在GitHub和Gitee上託管，兼顧國際化和國內速度。

\end{tcbinfo}

\textbf{進階玩法：用GitHub Pages搭建免費個人部落格}\par

GitHub Pages允許你免費託管一個靜態網站，配合Jekyll、Hexo、Hugo等靜態站點生成器，可以輕鬆搭建個人部落格。寫完Markdown文章push到倉庫，部落格自動更新，還自帶HTTPS和自定義域名。無數技術博主的部落格都是這麼搭的，零成本、無廣告、版本可控。

\booksection{九、總結：為什麼每個人都該瞭解GitHub}

GitHub早已不只是“程式設計師的工具”，它更像是一個\textbf{數字時代的工作方式}：版本管理讓一切可追溯，協作機制讓團隊高效運轉，開源精神讓知識自由流動。哪怕你不寫程式碼，理解“提交、分支、合併”這些概念，也能幫你更好地管理任何需要“反覆修改+多人協作”的內容——從論文到檔案，從筆記到方案。

對計算機系的學生來說，GitHub更是必備技能：它是你的作品集、協作平臺、學習資源庫，也是未來求職時最有說服力的“簡歷”。早點註冊帳號，把自己的課程作業、小專案傳上去，慢慢積累，幾年後你會擁有一份讓面試官眼前一亮的技術履歷。

\begin{tcbpractice}{來動手：完成你的GitHub初體驗}
\begin{enumerate}
  \item 註冊一個GitHub帳號，完善個人資料（頭像、簡介、感興趣的技術領域）。
  \item 建立一個名為\texttt{hello-\allowbreak github}的公開倉庫，勾選"Add a README file"。
  \item 進入倉庫，點選README.md的編輯按鈕（鉛筆圖示），用Markdown寫一段自我介紹，提交更改。
  \item 在GitHub搜尋欄搜尋一個你感興趣的關鍵字（比如\texttt{python-\allowbreak tutorial}），Star兩個你喜歡的專案。
  \item （進階）在電腦上安裝Git，把倉庫Clone到本地，修改README後push回雲端，完成一次完整的"本地雲端"工作流。
\end{enumerate}

\end{tcbpractice}
